\documentclass[a4paper,11pt]{article}
\usepackage[margin=27mm]{geometry}
\usepackage{fontspec,amsmath,amssymb,amsthm,unicode-math,booktabs}
\setmainfont{STIXGeneral}\setmathfont{STIX Two Math}
\newtheorem{theorem}{Theorem}\newtheorem{lemma}[theorem]{Lemma}
\newtheorem{proposition}[theorem]{Proposition}\newtheorem{corollary}[theorem]{Corollary}
\newcommand{\gen}[1]{\langle#1\rangle}\newcommand{\Aut}{\operatorname{Aut}}
\title{Faithfulness and the maps with at most two flag orbits}
\author{Human Compression project}\date{Working proof edition; independent review pending}
\begin{document}\maketitle
\begin{abstract}
We recover the principal faithfulness, symmetry and mixing results of
Cunningham, Moch\'an, Montero, Poto\v cnik and Toledo by applying the
common theory of equivalence relations and transporters. In rank three,
unfaithfulness reduces to one of two endpoint identifications. The remaining
classification is a calculation in a dihedral group, including infinite
polygons. The final construction in the source needs a change of one
generator in its displayed formula. Its converse also requires attention
to a degenerate choice of base generators. The corrected construction and
the finite example are stated separately from compression of valid results.
\end{abstract}

\section{Convention and source scope}
Write $I=\{0,\ldots,n-1\}$ and use the definitions in
\emph{Equivalence relations, transporters and flags}, cited below as the
common paper. Color words act on the right. Put
$F_i(x)=xW_{I\setminus\{i\}}$ and $C(x)=\bigcap_i F_i(x)$.
Thus $x\sim y$ means that $x$ and $y$ determine the same maximal chain
of faces. The symbols for involutions below refer to actions on flags;
when an automorphism is prescribed by its value at a base flag, it is
distinguished from the color connection having that value.

The source, abbreviated [F], is \emph{Faithfulness in maniplexes}, Journal
of Algebraic Combinatorics 64 (2026), article 23,
doi:10.1007/s10801-026-01559-y. Our exact edition has 29 physical pages.
The principal contract comprises its theory of face equivalence, its
classification of unfaithful maps with at most two flag orbits, and its
faithfulness and thinness theorems for mixes. Auxiliary pictured examples
are not themselves a classification census, and their individual
presentations are not regenerated here. Source errors are identified
explicitly rather than attributed to new theorems of the common paper.
The three final problems are recorded separately from this principal
contract.

\section{Color support and endpoint identifications}
\begin{lemma}\label{support}
Let $u$ and $v$ be color words with disjoint supports. If $x\sim xuv$,
then $x\sim xu$. If the two words commute, then also $x\sim xv$.
\end{lemma}
\begin{proof}
For a color $i$ outside the support of $u$, its path already shows
$xE_i xu$. For a color in that support, $v$ avoids $i$, so
$xuE_i xuvE_i x$. This accounts for every color. When the words
commute, interchange them and repeat the argument.
\end{proof}

Two $i$-adjacent flags have every other face in common. They are
chain-equivalent precisely when they have their $i$-face in common
as well. There is only one $i$-face precisely when this happens at
every flag: if every $i$-edge stays in an $i$-face, no edge leaves
that face, so connectedness makes it the whole graph. These facts
give [F, Lemma 3.1 and Propositions 3.2--3.3], including the
premaniplex case. A single face implies unfaithfulness whenever some
edge of its color is not a semiedge; no such implication is valid
for the one-flag premaniplex.

\begin{lemma}\label{endpoints}
If $x\sim y$, then $xr_0\sim yr_0$ and
$xr_{n-1}\sim yr_{n-1}$. If the maniplex is unfaithful, then for
each $1\leq i\leq n-2$ there are distinct equivalent flags joined
entirely by colors below $i$ or entirely by colors above $i$.
\end{lemma}
\begin{proof}
For the second assertion take distinct $x\sim y$. A path avoiding
$i$ can be written $uv$, where $u$ uses colors below $i$ and $v$
uses colors above $i$. The two words commute. Lemma~\ref{support}
puts both $xu$ and $xv$ in $C(x)$, and at least one differs from $x$,
since $xuv=y\ne x$.

For the first assertion every face of rank other than zero already
contains both $xr_0$ and $yr_0$. In particular a path between them
avoiding color 1 has the form $r_0^\epsilon w$, where $w$ uses only
colors $2,\ldots,n-1$ and $\epsilon\in\{0,1\}$. If $\epsilon=0$,
this is also a path avoiding zero, which proves the claim. Otherwise
$yr_0=xw$. Since $xE_0y$, this gives $yE_0yr_0$. Multiplying the
displayed identity on the right by $r_0$, and using commutation,
also gives $y=xr_0w$, hence $xE_0xr_0$. Transitivity now gives
$xr_0E_0yr_0$. Duality proves the other endpoint assertion.
\end{proof}

In rank three, the middle face consists of the four flags
$x,xr_0,xr_2,xr_0r_2$. Lemma~\ref{support} therefore shows that a
map is unfaithful exactly when some flag is equivalent to its
0-neighbor or its 2-neighbor. If $x\sim xr_1$, simplicity forces
$xr_1=xr_0r_2$; the same lemma then puts all four flags in $C(x)$.
This proves [F, Propositions 3.4--3.9].

The rest of [F, Section 3] also follows directly from face relations.
Parallel edges force their endpoints to be equivalent, whereas a path
with possible semiedges is faithful: its unique path between distinct
vertices cannot avoid every color. This gives the stated two- and
three-flag premaniplex tests. Covers preserve equivalence by mapping
color-avoiding paths to such paths. Equivalence inside a facet implies
equivalence in the ambient graph, since its witnessing paths also avoid
the new color. Thus ambient classes are unions of facet classes.
If every class has size at most two, the same holds in each facet.
If there is just one facet, and two distinct flags are equivalent
\emph{inside that facet}, they must be joined by the new-color edge:
both the old pair and that edge lie in an ambient class of size at most
two. This is the needed hypothesis in [F, Proposition 3.21]; merely
being joined by a path in the facet does not imply facet equivalence.
Finally two unique faces of different ranks give at least three distinct
equivalent flags $x,xr_i,xr_j$, ruling out classes of size at most two.

\section{Symmetry criteria without an orbit bound}
Let $G\leq\Aut(M)$. The common transport theorem applies to its free
action on the flags. Here $H_i(x)$ is precisely the setwise stabilizer
of $F_i(x)$. Its diagonal case gives [F, Theorem 4.2 and Corollary 4.3].
If $y\sim x$, it gives
\[
 yg\sim x\quad\Longleftrightarrow\quad g\in H(x),
 \qquad H(x)=\bigcap_iH_i(x).
\]
This is [F, Theorem 4.6]; its disjoint union of torsors gives
Corollaries 4.7--4.8, including the bound by the number of flag orbits
when $H(x)=1$. For different orbits, the full transporter formula and
quotient-path lift in the common paper give exactly Theorem 4.15 and
Corollary 4.16. No appeal to a finite enumeration is necessary.

Suppose now that there are two flag orbits, with $I_0$ the set of
colors preserving each orbit. Let $i,j\notin I_0$ be distinct.
For a base flag $x$, define $\tau_{ji}$ by
$x\tau_{ji}=xr_jr_i$. An automorphism with this value exists and
is unique, since the two color changes return to the base orbit.
The transporters from $xr_i$ to $x$ are $H_k(x)$ for $k\ne i$
and $\tau_{ji}H_i(x)$ for $k=i$. Indeed, deleting a color other
than $i$ permits the $i$-edge; for color $i$, the path using $j$
supplies the stated representative. Consequently
\[
 xr_i g\sim x\quad\Longleftrightarrow\quad
 g\in\tau_{ji}H_i(x)\cap\bigcap_{k\ne i}H_k(x).
 \tag{1}
\]
This is [F, Proposition 4.11]. Automorphism equivariance makes class
sizes constant on each orbit. Lemma~\ref{endpoints} makes them equal
across the two orbits whenever an endpoint color changes orbit,
proving Proposition 4.9.

For an orientable rotary maniplex, apply (1) to its rotation group
and $i=0,j=1$. If the distinguished rotations are
$\sigma_i$, with $x\sigma_i=xr_ir_{i-1}$, the criterion is
\[
 H(x)=1,\qquad
 S:=\sigma_1H_0(x)\cap\bigcap_{k=1}^{n-1}H_k(x)=\varnothing.
 \tag{2}
\]
This includes orientable regular maniplexes by using their index-two
rotation subgroup. If $\gen{\sigma_i}\cap\gen{\sigma_j}$ has a
nonidentity element for $|i-j|\geq2$, its two word representatives
together avoid every color, placing it in $H(x)$ and disproving
faithfulness. Finally $1\in S$ says $x\sim xr_0$, and
$\sigma_1\cdots\sigma_{n-1}\in S$ says
$x\sim xr_{n-1}$, because $xr_0\sigma_1\cdots\sigma_{n-1}
=xr_{n-1}$. Rotary symmetry and the corresponding endpoint adjacency
propagate each equality to every flag. Thus these conditions say
respectively that there is a single vertex or facet. This proves
[F, Corollaries 4.12--4.13 and Proposition 4.14] in the orientable
case. For a nonorientable regular maniplex, Proposition 4.14 alone
follows directly by $H_i=\gen{\rho_j:j\ne i}$ and
$\sigma_i=\rho_{i-1}\rho_i$. In that case the rotation group is the
whole group, so the orientable criterion (2) is not asserted. In particular
$1\in\sigma_1H_0$ is equivalent to
$\rho_0\in H_0$, and
$\sigma_1\cdots\sigma_{n-1}=\rho_0\rho_{n-1}\in H_{n-1}$
is equivalent to $\rho_{n-1}\in H_{n-1}$; the other memberships
are automatic when the corresponding endpoint subgroup is the whole
group.

\section{One-face maps and dihedral centralizers}
For a regular unfaithful map, an endpoint identification at one flag
holds at every flag. Up to duality there is therefore one facet, so
the flag set is the regular polygon action of
$D=\gen{a,b}$, where $a^2=b^2=1$. Let $p$ be the order of $ab$;
it is the number of edges of the facet, and may initially be infinite.
The third mark $c$ commutes with $a$ and differs from $1,a,b$.
The centralizer of a reflection in a finite dihedral group is
$\{1,a\}$ when $p$ is odd and $\{1,a,h,ha\}$ when $p$ is even,
where $h=(ab)^{p/2}$; in the infinite group it is $\{1,a\}$.
These assertions follow by writing an element as $(ab)^m$ or
$(ab)^ma$ and using $a(ab)a=(ab)^{-1}$.
Thus $p=2k$ is finite and $c=h$ or $ha$. These are precisely the
regular maps called $P\{k\}$ and $T\{k\}$ in [F, Proposition 5.1],
and their duals. Parameters producing repeated marks are omitted:
in particular $T\{1\}$ is not a maniplex. The underlying group has
order $4k$; taking the orders of $bc$ gives the source's types.
For $c=h$, this order is 2. For $c=ha$, it is $k$ when $k$ is odd
and $2k$ when $k$ is even.

Now suppose there are exactly two flag orbits. If a color endpoint
changes orbit, its own edge or the other endpoint edge propagates
an endpoint identification to both orbits, by
Lemma~\ref{endpoints}. Hence an unfaithful two-orbit map has a
single vertex or facet unless its symmetry type is $2_{\{0,2\}}$.
This proves [F, Proposition 5.2].

The one-facet case can be treated without a separate presentation
for every symmetry type. Its automorphism group is an index-two
subgroup $H$ of the polygon group $D$: the full polygon group acts
regularly on flags, and $H$ has two orbits. Write its character
as $\chi:D\to C_2$. A color connection commutes with the left
$H$-action and is determined by its values at one flag in each orbit.
Let $a,b$ again be the polygon's color 0 and color 1 connections.

If $\chi(a)=\chi(b)=1$, then $H$ is the rotation group. Let $w$
be the third neighbor of the identity flag. Commutation with $a$
determines the connection on the second orbit. If $w\in H$,
involutivity forces $w^2=1$; if $w\notin H$, it forces $(wa)^2=1$.
In either case the only allowed nontrivial rotation is the central
half-turn. The third connection is then one global right multiplication,
so the full polygon group acts and the map is regular. This excludes
the two-orbit types $2_\varnothing$ and $2_{\{2\}}$ with one facet.
If $\chi(a)=0,\chi(b)=1$ and the third color changes orbit, its
value $w\notin H$ must commute with $a$. The same centralizer
calculation gives $w=h$ or $ha$. Involutivity then makes the
connection global right multiplication, again giving regularity.
This excludes type $2_{\{0\}}$ with one facet. Duality handles
the corresponding one-vertex cases. Together with the previous
paragraph this proves [F, Propositions 5.3--5.4].

\section{The other one-facet types}
It remains to take $\chi(a)=1,\chi(b)=0$. The polygon must have
$p=2k$ edges, or be infinite; color its vertices alternately black
and white. Its two $H$-orbits are those two flag colors. Choose
an involution $w\in H$, and define the third connection by
\[
 R_2(x)=
 \begin{cases}xw,&x\text{ black},\\ xawa,&x\text{ white}.
 \end{cases}
 \tag{3}
\]
The first two connections are $a,b$. The rule commutes with $a$
because that connection interchanges the colors. It is an involution
because $w$ preserves color. It is fixed-point free exactly when
$w\ne1$, never coincides with $a$, and coincides with $b$ precisely
when $w=b$ or $aba$. All these conclusions use the free polygon
action and therefore hold for infinite polygons as well.

Every color-preserving polygon automorphism commutes with (3).
A color-changing one can exist precisely when the automorphism
mapping the base flag to its 0-neighbor exists: the preserving
subgroup already acts transitively on each color. This condition
is $aw=wa$. In the finite case its additional possibilities are
$h=(ba)^k$ and $ah$; in the infinite case there are none.
Thus (3) gives exactly the unfaithful maps of type $2_{\{1,2\}}$
when
\[
 w\notin\{1,b,aba,h,ah\},\qquad w\in\ker\chi,
 \quad w^2=1.
 \tag{4}
\]
The terms $h,ah$ are omitted for an infinite polygon. Conversely
any map of this type has one facet: it has two vertex orbits, so
cannot have a single vertex. Equivariance determines its third
connection on the black flags by one element $w$; commutation with
$a$ determines it on the white flags. This proves (3)--(4), and
hence [F, Proposition 5.6].

Duality gives type $2_{\{0,1\}}$. Replacing color 2 by the
product of colors 0 and 2 is the opposite operation. Applied to
(3), it gives type $2_{\{1\}}$ with one facet. Its third
connection now changes color, so cannot coincide with $b$;
the exclusions $b,aba$ in (4) are no longer needed. The exclusion
of 1 prevents coincidence with $a$, and $h,ah$ still exclude
regularity. Duality gives the one-vertex case, equivalently the
Petrie dual of the dual of (3). These operations are involutive,
so their converses follow by reversing them. This proves all four
parts of [F, Theorem 5.8], with the infinite-parameter convention
made explicit. The literal occurrence of $a$ in [F, Remark 5.7]
as a color-preserving exceptional choice is inapplicable:
$\chi(a)=1$; the possible repeated-edge choice there is $b$.

\section{The remaining symmetry type}
The complete construction and proof for type $2_{\{0,2\}}$ are
given in the companion \emph{A corrected construction for unfaithful maps
with two flag orbits}.
Its input is a regular map with finite facets of even length.
In coordinates $G\times\{0,1\}$, let its marks be $a,b,c$,
let $p=|ab|$ be even, and put $h=(ab)^{p/2}$. The new marks are
\[
 R_0(g,0)=(ga,0),\quad R_0(g,1)=(gb,1),\quad
 R_1(g,e)=(g,1-e),
\]
\[
 R_2(g,0)=(gc,0),\qquad
 R_2(g,1)=(gd,1),\quad d=h\text{ or }hb.
 \tag{5}
\]
These constructions and their duals exhaust the unfaithful maps
of the remaining type, with the regular cases removed. The source's
printed torus rule uses $ha$ in place of $hb$. For $p>2$, $ha$
does not commute with the white color-0 mark $b$; for $p=2$ it
equals $b$, giving parallel edges. Thus that literal rule never
gives a maniplex on a valid polygon facet. The companion proves
the corrected criterion and explains the swap of the two flag
orbits needed when an initial base marking repeats a generator.
This is a correction to the source's formulation, not a claimed
solution of one of its open problems.
The following unnumbered source paragraph about the exact sizes of all
chain-equivalence classes is outside this recovery: the construction's
validity and exhaustive converse do not by themselves prove that assertion.

\section{Covers, mixes and thinness}
A cover of a faithful maniplex that is injective on each facet is
faithful by the second part of the common paper's joint-map theorem.
This is [F, Proposition 6.1]. The coordinate projections of a mix
jointly distinguish its flags and preserve face equivalence. Thus
\[
 C_{X\diamond Y}(x,y)\subseteq C_X(x)\times C_Y(y).
\]
The coordinate pairs in the first sentence of [F, Proposition 6.2]
are transposed: the intended implications are $x\sim x'$ and
$y\sim y'$, for equivalent pairs $(x,y),(x',y')$. The displayed
inclusion and subsequent corollary in the source have the correct
meaning. The common theorem now gives its Corollaries 6.3--6.4,
and even removes the finite-orbit restriction from the latter.

The common paper's deleted-face criterion is precisely the
equivalence in [F, Propositions 7.1--7.2]. Applying it to the mix,
the two-coloring forces both coordinates to choose the same flag
or the same adjacent flag. This gives [F, Theorem 7.3 and Corollary
7.4]. No claim is made that orientability can be dropped. The
source's cited nonorientable example remains an attributed example,
not an additional finite census carried out here.

\section{What is recovered and what remains separate}
The generic transporter theorem recovers all the symmetry criteria
without a bound on the number of orbits. The support lemma and
endpoint propagation reduce the rank-three classification to the
centralizer of a reflection. The common paper also proves all
covering and mixing implications used here. These are complete
written arguments; they do not rely on sampled finite groups.
The companion supplies the last construction with its corrected
literal formula and full converse.

The attribution remains with the authors of [F] for the original
results, and with the earlier authors cited there for the rank-three
endpoint criterion and classical one-face maps. The new organizational
step is to use one common refinement and transporter calculation
throughout. Standard facts used in that calculation are proved in
the common paper. The additional arbitrary-family result is not
assigned historical novelty. The reconstruction criterion is kept
separate from recovery of prior results: it treats natural
reconstruction in Problem 3, and arbitrary isomorphism for finite
posets, with the remaining interpretation stated explicitly.
Problems 1 and 2 are not claimed solved. Nor does the present
application claim complete Lean formalization of the classification.
\end{document}
